\section*{Erd\H{o}s Problem 799: list colouring a typical graph}

\paragraph{Statement and claim.}
For \(G\sim G(n,1/2)\), with probability tending to one,
\[
 \chi_L(G)=O\left(\frac{n\log\log n}{\log n}\right)=o(n).
\]
We prove this using a random allocation of list colours
and a reserve palette. The external graph-theoretic input
is the classical ordinary-colouring bound
\(\chi(G(n,1/2))=O(n/\log n)\) with probability tending to one.
This bound, and its use in the known solution, are recorded
in Section 4 of Alon's primary paper
\emph{Choice numbers of graphs: a probabilistic approach},
\url{https://www.cs.tau.ac.il/~nogaa/PDFS/choice.pdf}.
See \url{https://www.erdosproblems.com/799}.

\paragraph{A typical graph has small colour classes.}
Let \(s=\lceil3\log_2 n\rceil\). By the first-moment bound,
\[
 \Pr(\alpha(G)\geq s)
 \leq\binom ns2^{-\binom s2}\longrightarrow0.
\]
Together with the stated ordinary-colouring theorem,
we may therefore fix a graph having a proper partition
into \(r\leq K n/\log n\) independent classes, each of size
at most \(s\). Necessarily \(r\geq n/s\).
We will prove a list-colouring bound for every such graph,
so there will be no union bound over assignments of lists.

\paragraph{Allocating colours, including a reserve.}
Fix arbitrary lists of
\(k=\lceil16r\log s\rceil\) distinct colours at every vertex.
Independently for each colour appearing anywhere in the
lists, assign it to a reserve palette with probability \(1/2\),
and to each of the \(r\) class palettes with probability \(1/(2r)\).
These are disjoint palettes of colour names; the same colour
has the same allocation at every vertex.
Call a vertex bad if its list misses the palette belonging
to its proper colour class.
For every fixed vertex,
\[
 \Pr(\text{bad})=(1-1/(2r))^k\leq e^{-k/(2r)}
 \leq s^{-8}.
\]
Thus the expected number of bad vertices is at most
\(ns^{-8}\leq rs^{-7}\), and Markov's inequality gives
\(\Pr(\text{at least }r\text{ bad vertices})\leq s^{-7}\).

The number of reserve colours in any one list is
\(\operatorname{Bin}(k,1/2)\).
The elementary Chernoff bound gives
\(\Pr(\text{fewer than }k/4\text{ reserve colours})\leq e^{-k/16}\).
By the union bound, the probability of any such deficient
list is at most \(ne^{-k/16}=o(1)\), since
\(k\geq16(n/s)\log s\) and \(s=O(\log n)\).
For sufficiently large \(n\), both desirable properties
therefore hold simultaneously for some allocation:
fewer than \(r\) bad vertices, and at least \(k/4\geq r\)
reserve colours in every list.

\paragraph{Completing the list colouring.}
Every good vertex chooses a colour from its own class palette.
Adjacent good vertices lie in different classes and use
disjoint palettes, so this is proper.
Colour the fewer than \(r\) bad vertices one at a time,
using different reserve colours.
Each has at least \(r\) available reserve colours in its list,
and fewer than \(r\) have previously been used.
These colours cannot conflict with good vertices because
the reserve is disjoint from all class palettes.
This gives a proper colouring from the originally arbitrary
lists, proving \(\chi_L(G)\leq k\).

\paragraph{Conclusion and scope.}
Substituting \(r=O(n/\log n)\) and \(s=O(\log n)\) gives
the claimed \(o(n)\) bound.
The probability concerns the random graph; for every
graph in the typical class, the deterministic conclusion
holds for every assignment of lists.
The sharper known order \(\Theta(n/\log n)\) is due to
Alon--Krivelevich--Sudakov; it is not needed for the
complete affirmative answer proved here.
The only substantial external input in this reconstruction
is the ordinary-colouring theorem, not a list-colouring theorem.

\paragraph{Completion estimate.}
\textbf{COMPLETION: 100\%}
